\odiseachapter{Introducción}{De libros y películas}\label{cap:intro}

En julio de 1982, yo acababa de terminar el cuarto curso de la licenciatura en Ciencias Físicas y había ganado una beca para hacer prácticas de verano en la central nuclear de Cofrentes, por entonces todavía en construcción. Las prácticas empezaban en unos pocos días y aquella mañana regresaba de las oficinas de Hidroeléctrica (hoy Iberdrola) en Valencia, donde había firmado el primer contrato de mi vida laboral. De camino a mi piso de estudiantes, me detuve en una de las librerías universitarias cercanas a Blasco Ibáñez, con el ánimo de hacerme con una buena novela para entretenerme durante los viajes entre Requena (donde estaba previsto que me alojara) y Cofrentes. Entre ambos pueblos había una hora de carretera. No teníamos redes sociales ni teléfonos móviles. Yo leía mucho, es cierto, pero leer no era nada tan extraño en los ochenta. El autobús que tomaría cada día durante los siguientes dos meses iba lleno de trabajadores de todos los oficios: electricistas, mecánicos, oficinistas, ingenieros y estudiantes en prácticas como yo. Y casi todos llevábamos un libro a cuestas para el trayecto.

El que me asaltó desde el escaparate de la librería aquella mañana fue nada menos que \emph{El señor de los anillos}. Eran tres volúmenes, editados por Minotauro, y no eran baratos, pero los compré sin vacilar a cuenta de mi primer salario. Había oído hablar de Tolkien, porque alguien había contado una historia de las revueltas estudiantiles del mayo del 68 en París, cuando en el metro aparecían pintadas que afirmaban en rojo vivo sobre las paredes «Frodo vive». Yo no sabía muy bien quién era Frodo ni qué tenía que ver con la revolución, pero cualquier sombra de duda se despejó apenas leí los versos con los que arrancaba la trilogía.

\begin{versocita}
Tres anillos para los reyes elfos bajo el cielo.\\
Siete para los señores enanos en palacios de piedra.\\
Nueve para los hombres mortales condenados a morir.\\
Uno para el señor oscuro, sobre el trono oscuro\\
en la tierra de Mordor donde se extienden las sombras.
\end{versocita}

¿Por qué aquellos versos fueron suficientes para vaciarme los bolsillos y comprometerme con mil páginas de relato épico? Creo que la respuesta tiene que ver con el hecho de que mi cabeza estaba llena de palabras, leídas en las innumerables novelas y libros de poesía que ya había devorado por entonces. Palabras que resonaron con aquellas otras palabras del poema inicial, como un diapasón resuena al sonido de otro.

La trilogía de Peter Jackson llegó veinte años más tarde. Naturalmente, hice cola en el cine para verla y, naturalmente, la he vuelto a ver muchas veces desde entonces. La vi cuando mi hija Irene era un bebé y mi hijo Héctor aún no había nacido, la vi otra vez cuando aún eran tan niños que los abuelos nos reñían por darles sustos con los Nazgûl, y otra vez cuando eran un poco más mayores, y otra vez de adolescentes y otra vez durante la covid y otra vez recientemente. Por culpa del pesado de su padre, mis hijos se saben la trilogía de memoria.

Y sin embargo, ninguno de los dos ha leído los libros.

Los libros han estado siempre en casa: la versión española que compré en 1982 y la versión en inglés que compré en una librería de viejo en San Francisco hacia 1989, durante mi posdoctorado en California. Están muy visibles en las estanterías, provocadores y sensuales como las sirenas, con su excelente encuadernación y sus mil páginas de prodigios. En la misma estantería están \emph{Las mil y una noches}, los poemas completos de Czesław Miłosz y los de Raymond Carver, \emph{El Aleph} de Borges, \emph{Rayuela} de Cortázar, \emph{Ana Karenina} de Tolstoi y \emph{Los desposeídos} de Ursula K. Le Guin. Todos esos volúmenes desean la misma atención que esperaban los libros de la biblioteca de mi padre. Pero yo, a aquellos, no los decepcioné. Los de la mía esperan en vano a Godot.

Irene es médica; Héctor está acabando su máster en Ciencias Físicas. Los dos leían de pequeños. Son de la generación de Harry Potter y además en casa no teníamos televisión, así que, durante su niñez, como para su padre antes que ellos, los libros eran un excelente antídoto contra el aburrimiento. Ambos siguen leyendo hoy en día. Irene tiene poco tiempo, pero lo compensa con un extraño sexto sentido para seleccionar novelas interesantes. De hecho, fue ella la que me descubrió a una de mis escritoras favoritas, Elena Ferrante, y leímos juntos, gozosamente, la tetralogía de \emph{La amiga estupenda}. Héctor acaba de leer \emph{La insoportable levedad del ser} y los dos hemos pasado muchas horas hablando de la honda, casi inaudita percepción del alma humana que allí ofrece Milan Kundera.

¿Entonces? La diferencia entre mis hijos y yo no es de actitud, ni de interés. La diferencia está en los números. Aunque a ellos las redes sociales les sorprendieron relativamente tarde, no han sido inmunes a un asalto que ha mermado el tiempo de todo el mundo, sobre todo de la gente joven, robándoselo a muchas otras actividades más enriquecedoras para el alma, en particular la lectura. En el mundo en el que yo crecí, la televisión ofrecía una programación fija durante unas pocas horas del día. Empezaba a las seis o las siete de la tarde y terminaba dos horas después de la cena. No había internet y el cine era un lujo reservado a los domingos. No había, por supuesto, TikTok ni Instagram. Yo leía antes de dormir, leía los fines de semana, en vacaciones, durante los viajes, a cualquier hora. No lo hacía por devoción intelectual o por adquirir cultura. A menudo, no había nada mejor que hacer. Pero una infancia y una adolescencia de lecturas a destajo construyen un vocabulario y una imaginería a los que uno ya no puede sustraerse.

Los elfos de \emph{El señor de los anillos}, los enanos y los orcos, los Rohirrim, las temibles batallas del Abismo de Helm y el viaje final de Frodo por Mordor se formaron en mi mente antes de ver una sola imagen en la pantalla. Para mis hijos y para casi todos los jóvenes del planeta, todos esos seres fantásticos, todas esas tierras misteriosas son las que Peter Jackson imaginó. En tiempos de libros, había tantos elfos como lectores. Hoy en día, Galadriel tiene el rostro de Cate Blanchett, Arwen es Liv Tyler, Frodo es Elijah Wood. Nada que objetar, todos ellos son magníficas encarnaciones de los personajes de la leyenda. Pero cuando el verbo se hace carne, irremediablemente, pierde parte de su carácter divino.

Mis hijos, como muchos jóvenes en la franja entre los veinte y los treinta años, han leído menos que sus padres, pero la lectura ha estado y sigue estando presente en sus vidas. Entre los de menor edad, sin embargo, es frecuente que lean ---odiosos deberes aparte--- solo los \emph{feeds} de sus redes sociales. Peor aún, es un secreto a voces que la capacidad de comprender un texto complejo se ha desplomado entre los jóvenes. No menos sabido es el hecho de que su vocabulario se ha reducido. Y si combinamos ambas cosas, el resultado inevitable es una generación con menos recursos expresivos, menos capacidad analítica, menos instrumentos de raciocinio. O, en otras palabras, una generación de jóvenes con menos capacidad para pensar por sí mismos, por tanto, quizás más crédulos, menos críticos, más fáciles de manipular, más susceptibles a las estafas del populismo, las teorías de la conspiración o la demoledora tontería de los \emph{influencers}.

Ahora bien, como afirma Rick Blaine ---tan indistinguible de Humphrey Bogart como Aragorn de Viggo Mortensen---, siempre nos quedará París, es decir, el cine. Cierto, mis hijos no han leído \emph{El señor de los anillos}, pero podría argumentarse que las películas son tan buenas que una cosa compensa la otra. Y quizás lo mismo ocurre con todas las grandes historias. No hay gran libro, se diría, que no cuente con su versión cinematográfica y, en algunos casos, el celuloide supera de largo a la página escrita. \emph{Blade Runner}, por ejemplo, le da vuelta y media a la novela original (\emph{Do Androids Dream of Electric Sheep?}) y la trilogía de \emph{El padrino} reside en una esfera del Parnaso a la que no habrían llegado jamás las novelas de Mario Puzo. \emph{Muerte en Venecia} (Mann) o \emph{Los muertos} (Joyce) son obras maestras en el papel y en el cine. La serie televisiva sobre \emph{La amiga estupenda} de Ferrante es tan buena como el texto escrito, entre otras cosas porque la autora participó activamente en el guion. La lista es larga y habría mucho que escribir sobre la intertextualidad del cine y la novela, sobre la miríada de formas en que se complementan y se completan.

Entonces, ¿podemos prescindir alegremente del anticuado papel que nos obliga a descifrar signos y construir imágenes en el aire? ¿Podemos fiarle nuestra cultura, es decir, nuestra formación como ciudadanos, al cine, a las series televisivas? ¿Se pierde algo esencial por no haber leído a Tolkien? Quizás los obsesos de la lectura son (somos) \emph{boomers} nostálgicos que no se enteran de que los tiempos han cambiado, como esos programadores chapados a la antigua que aún no han admitido que todo su conocimiento de C++, o de cualquier otro lenguaje de programación, es inútil en los tiempos de la IA.

¡Ay!, mucho me temo que cuando el cine ---o la serie televisiva--- reemplaza por completo a la lectura, hay algo que se queja. Cine y literatura tienen su espacio propio, sus recursos únicos ---también sus limitaciones---, y en mi opinión, forman una pareja perfecta, uno de esos matrimonios enamorados que son como un magnífico animal bicéfalo, cada parte de la pareja un individuo único, pero el total mayor que la suma de sus partes. Cuando la lectura sufre, el cine sufre. Durante tres mil años, los humanos han usado los libros para construir el mundo. Se diría, sin embargo, que vienen tiempos iletrados, tiempos donde no hará falta quemar los libros como en \emph{Fahrenheit 451} porque nos limitaremos a ignorarlos, tiempos donde toda la información y todo el entretenimiento van a ser audiovisuales, tiempos donde incluso el cine se va a ver amenazado. Porque, ¿qué joven tiene paciencia para sentarse quieto dos o tres horas delante de una pantalla, salvo para seguir las aventuras de un adolescente en pijama colgando de una tela de araña?

Por otra parte, una película taquillera tiene que gustarle a millones de personas, lo que implica, además de actores famosos y mucho espectáculo, un argumento que no ponga nervioso a nadie. Por ejemplo, un jovencito con poderes arácnidos, dando fabulosos saltos y combatiendo a malos suficientemente domesticados.

¿Qué ocurre entonces cuando un ambicioso director se propone narrar una de las grandes épicas de la literatura? Sabemos, por \emph{El señor de los anillos}, que es posible hacerlo, pero no olvidemos que Tolkien es tan grandioso como poco complicado. Los buenos carecen de fisuras; los malos son feísimos; el talante moral de Aragorn, de Arwen, de Gandalf o de Legolas es el mismo que el del Capitán Trueno; Galadriel ofrece una pizquita de oscuridad semejante a la aceituna en el martini de bondad de los protas. Frodo y Sam protagonizan un \emph{bromance} impecable; Boromir se hace el malo dos minutos y a continuación muere por la causa. El único personaje complejo es Gollum, que, por otra parte, acaba por redimirse y sacar a Frodo de apuros. Con toda mi admiración hacia Jackson, la trilogía de Tolkien se dejaba querer.

¿Y la \emph{Odisea}? En este caso ha habido varios intentos con diferentes méritos y limitaciones. El \emph{Ulises} de Camerini, estrenado en 1954 y protagonizado por Kirk Douglas, tenía escenas afortunadas ---incluyendo la del cíclope--- y trucos muy innovadores, como usar la misma actriz para encarnar a Penélope y Circe. Recuerdo haber visto esta película en el cine de Alba\footnote{De hecho, mi adicción al cine es casi tan temprana como mi vicio por la lectura. A los diez años, mi familia se mudó a un pueblo de la costa mediterránea, situado entre Vinaroz y Amposta, al que siempre me he referido como Alba, quizás porque, como Cervantes, prefiero no acordarme de su nombre. Mi padre era la autoridad militar de la plaza, lo que nos daba a sus hijos la identidad de extraterrestres y a mí me granjeó durante mucho tiempo el estatus de paria. Durante la semana, me veía a menudo con Quelo, mi único amigo, pero los domingos Quelo tenía otros planes, con sus amigos de siempre. Yo solo tenía la sesión doble en el cine del pueblo y una bolsa de pipas, para olvidar la soledad.} con doce o trece años y haberme pasado muchas tardes jugando a ser Ulises y Telémaco, a tensar el arco y asaetear a los pretendientes, empezando por Antínoo, que se parecía mucho a Adrián\footnote{Cuyo verdadero nombre, huelga decirlo, tampoco era ese.}, el abusón que me traía por la calle de la amargura en el instituto. Naturalmente, era demasiado joven para entretenerme en sutiles análisis de contenido, pero agradezco de todo corazón la compañía de Kirk Douglas en mis juegos solitarios.

Así que cuando el celebrado director Christopher Nolan anunció sus planes de rodar una nueva versión de la
\emph{Odisea} (a la que podríamos llamar, con un poco de descaro, la \emph{Nodisea}), pensé que se trataba de buenas noticias. Con Nolan pasa lo que con el torero Curro Romero, del que se decía que había que ir siempre a sus corridas \emph{por si acaso}.

Aunque esta corrida yo se la auguraba complicada. El problema es que Homero es mucho más difícil de manejar que Tolkien. Odiseo es un personaje casi inabarcable. Astuto, sagaz, mentiroso, artero, cruel, implacable, odioso, valiente, elegante, generoso a veces, buen hijo y mejor padre, mal marido, gran amante, observador, sensible, capaz de hacerse atar a un mástil para escuchar a las sirenas y de arrojar a un bebé desde las murallas de Troya.

Un personaje así presenta un problema serio para el cine, sobre todo para un director \emph{superstar} que quiere rodar una película taquillera. Pero quizás, me dije a mí mismo, sir Christopher demostraría estar a la altura de tamaño reto y nos ofrecería, si no una lectura literal ---imposible en tres horas de película---, al menos una interesante versión del héroe que enriqueciera el mito. La película se estrenó en julio de 2026. Unos pocos días después, fui a verla. Como todo el mundo,
he leído las mil y una críticas que se le han dedicado. Muchas de estas críticas alaban o denuestan las elecciones puramente cinematográficas del director. Se han vertido ríos de tinta sobre cuestiones tales como el color de la piel de Helena, la geometría del caballo o el diseño de las armaduras. Por mi parte no tengo gran cosa que decir al respecto. Nolan, como todo director que se precie, tiene sus virtudes, sus defectos, sus golpes de inspiración y sus manías, y en la película hay un poco de todo eso. El problema de la \emph{Nodisea}, en mi opinión, no tiene nada que ver con la construcción formal de la película sino con la versión del mito que el celebrado director y Hollywood nos ofrecen.

Casi a la vez que se estrenaba la \emph{Nodisea}, se publicaba la maravillosa y atrevida traducción de la \emph{Odisea} de Juan Manuel Rodríguez Tobal, bajo el prestigioso sello de Hiperión. Decir que su lectura me impactó mucho más que la película sería el equivalente a comparar la picadura del áspid con la del mosquito. De hecho, sin Tobal, este libro no se habría escrito. En el capítulo~\ref{cap:osada} explico con detalle los porqués, que, sin embargo, pueden resumirse en una sola línea. La obra consigue lo único que justifica atreverse a traducir a Homero en estos tiempos iletrados.

Es decir, hacer magia.

El objetivo de este libro es acercar a la querida lectora, al estimado lector, al Poeta con mayúsculas, al autor ---al menos según Borges--- de dos de las cuatro grandes narraciones que fundan la literatura occidental. Propongo aproximarnos a Homero mediante la lectura simultánea de la \emph{Odisea} ---de la mano de Tobal--- y dos versiones modernas de esta. La que nos ofrece Nolan ---es decir, Hollywood--- y la que aventuro en \emph{Abandonando Ítaca} (Eolas, 2025)\footnote{\emph{Abandonando Ítaca}, \url{https://eolasediciones.es/libro/abandonando-itaca-leaving-ithaca/}.}, un poema en el que imagino a un Ulises octogenario que repasa su vida en una Ítaca que ya no siente como suya. En el penúltimo capítulo de este libro, veremos una cuarta versión, la que imaginó Robert Graves en su magnífica novela \emph{La hija de Homero}, publicada en 1955.

El hecho de que me atreva a ofrecer tres versiones modernas de las aventuras de Odiseo, incluida la mía, se debe a mi convicción de que el mito existe para reinterpretarlo y cada época está en su derecho de hacerlo como le parezca conveniente. Ese no es el problema.
El problema, como veremos, es que el cine actual ---el cine comercial, destinado a entretener, que no a educar a las masas--- teme a Homero. Y lo teme, porque el Poeta nos pone un espejo delante en el que se refleja nuestra propia crueldad y falta de escrúpulos. Quizás por eso, el director más popular del momento ha convertido al héroe más amoral e inclasificable de la historia de la literatura en un cliché santurrón, reflejando nuestra infinita hipocresía.
